\section{Some background material
\label{back}
}

We gather here some necessary facts and def\/initions from complex analysis.
A good reference is~\cite{GR}.

We work in the category of complex spaces, i.e.\
${\mathbb C}$-ringed spaces locally modelled on ${\mathbb C}$-ringed subspaces of domains $U$ of ${\mathbb C}^n$ def\/ined
by f\/initely many holomorphic functions in $U$.
In particular, we allow nilpotents in the structure sheaf.

For a~complex space $(X,{\mathcal{O}}_X)$, the topological (Urysohn--Menger), analytic (Chevalley) and algebraic (Krull)
notions of local dimension $\dim_xX$ coincide.
A complex space is called {\em reduced} if its structure sheaf has no nilpotents.
It is {\em pure dimensional} (or equidimensional) if $\dim_x X$ is the same at all points of $X$.
For a~$0$-dimensional space $X$, its {\em length} is the sum of dimensions of stalks ${\mathcal{O}}_{X,x}$.

The {\em cotangent sheaf}\, $\Omega^1_X$ of a~complex space $X$ is def\/ined by glueing together the cotangent sheaves
${\mathcal{I}}_V/{\mathcal{I}}_V^2$ of local model spaces $(V, {\mathcal{O}}_{{\mathbb C}^n}/{\mathcal{I}}_V)$.
The {\em tangent sheaf}\, ${\mathcal{T}}_X$ is its dual $\operatorname{Hom}(\Omega^1_X,{\mathcal{O}}_X)$.
The {\em tangent space}\, at $x\in X$ is $T_xX={\mathcal{T}}_{X,x}/\mathfrak{m}_x {\mathcal{T}}_{X,x}$, where
$\mathfrak{m}_x$ is the maximal ideal of the stalk ${\mathcal{O}}_{X,x}$.

If $A$ is a~local ring with maximal ideal $\mathfrak{m}$ and $M$ is a~module over $A$, then a~sequence
$(f_1,\dots,f_r)\in \mathfrak{m}$ is said to be {\em regular for $M$} if $f_{i+1}$ does not divide $0$ in
$M/(f_1M+\dots+f_iM)$ for $i=0,\dots, r-1$.
The length of a~maximal regular sequence is called the {\em depth of $M$}.
A~complex space $X$ is said to be {\em Cohen--Macaulay} if $\operatorname{depth}{\mathcal{O}}_{X,x}=\dim_xX$ for all
$x\in X$.
A~reduced $1$-dimensional complex space is always Cohen--Macaulay.

A closed complex subspace $X$ of a~complex space $Z$ is called {\em regularly embedded} of codimension~$r$ if every of
its local def\/ining ideals $ {\mathcal{I}}_{X,x}$ has depth~$r$.
If $Z$ is a~complex manifold, and $X\subset Z$ is regularly embedded, then~$X$ is called a~{\em local complete
intersection}, abbreviated l.c.i.
This is an intrinsic condition and it implies that $X$ is Cohen--Macaulay.

For a~closed complex subspace $X$ of a~complex space $Z$ with ideal sheaf $ {\mathcal{I}}_X\subset {\mathcal{O}}_Z$, the
normal sheaf of $X$ in $Z$ is ${\mathcal{N}}_{X/Z}=\operatorname{Hom}_{{\mathcal{O}}_X}( {\mathcal{I}}_X/
{\mathcal{I}}_X^2,{\mathcal{O}}_X)$.
If $X\subset Z$ is regularly embedded, then both ${\mathcal{N}}_{X/Z}$ and $ {\mathcal{I}}_X/ {\mathcal{I}}_X^2$ are
locally free.
If, moreover, $Z$ is smooth and $X$ is reduced, then the conormal sheaf $ {\mathcal{I}}_X/ {\mathcal{I}}_X^2$ f\/its into
an exact sequence
\begin{gather*}
0\to {\mathcal{I}}_X/ {\mathcal{I}}_X^2\to \Omega^1_Z|_X\to \Omega^1_X\to 0,
\end{gather*}
and consequently we have an exact sequence
\begin{gather*}
0\to {\mathcal{T}}_X\to {\mathcal{T}}_Z|_X\to {\mathcal{N}}_{X/Z}\to \mathcal{E}xt^1\big(\Omega^1_X,{\mathcal{O}}_X\big).
\end{gather*}

A celebrated theorem of Douady~\cite{Dou} states that for any complex space $Z$ there exists a~complex space
${\mathcal{D}}(Z)$ parameterising all pure-dimensional compact complex subspaces of $Z$.
In addition, there exists a~{\em universal family} on ${\mathcal{D}}(Z)$, i.e.\
a~complex subspace $Y$ of ${\mathcal{D}}(Z)\times Z$, which def\/ines a~double f\/ibration
\begin{gather}
{\mathcal{D}}(Z)\stackrel{\nu}{\longleftarrow} Y \stackrel{\mu}{\longrightarrow} Z
\label{seqq0}
\end{gather}
with the following properties:
\begin{itemize}\itemsep=0pt
\item[(i)] $\nu$ is f\/lat and proper; \item[(ii)] if $S$ is a~complex space, $T\subset S\times Z$ a~complex subspace with
properties stated in (i), then there exists a~unique holomorphic map $f:S\to {\mathcal{D}}(Z)$ such that $T\simeq
S\times_{{\mathcal{D}}(Z)} Y$.
\end{itemize}
The f\/ibre of $Y$ over any $m\in {\mathcal{D}}(Z)$ is the complex subspace $X$ of $Z$ corresponding to $m$ and the
restriction of the normal sheaf ${\mathcal{N}}_Y$ of $Y$ in ${\mathcal{D}}(Z)\times Z$ to this f\/ibre is
${\mathcal{N}}_{X/Z}$.

At any $X\in {\mathcal{D}}(Z)$, there is a~canonical isomorphism $T_X{\mathcal{D}}(Z)\simeq H^0(X,{\mathcal{N}}_{X/Z})$.
The Douady space ${\mathcal{D}}(Z)$ is smooth at $X$ if $\operatorname{Ext}^1_{{\mathcal{O}}_Z}(
{\mathcal{I}}_X,{\mathcal{O}}_X)=0$.
For a~regularly embedded $X$, this is equivalent to $H^1(X,{\mathcal{N}}_{X/Z})=0$.
For a~projective $Z$, ${\mathcal{D}}(Z)$ is the Hilbert scheme parameterising compact subschemes of $Z$.

\section{Hypercomplex and hyperk\"ahler structures\\ from higher degree curves}\label{hgc}

For us, a~{\em curve} means a~compact, pure $1$-dimensional, Cohen--Macaulay complex space.
If the curve is reduced, then the Cohen--Macaulay assumption is redundant.

Let $Z$ be a~pure $(n+1)$-dimensional complex space and $\pi:Z\to {\mathbb{P}}^1$ a~holomorphic surjection.
For a~closed subspace $X$ of $Z$ we continue to write $\pi$ for its restriction to $X$.
Each f\/ibre $X_\zeta$, i.e.\
the complex subspace $(\pi^{-1}(\zeta), {\mathcal{O}}_X/\pi^\ast\mathfrak{m}_\zeta$), is a~closed subspace of~$X$.
If ${\mathcal{F}}$ is a~sheaf on $X$, we write ${\mathcal{F}}(k)$ for ${\mathcal{F}}\otimes \pi^\ast
{\mathcal{O}}_{{\mathbb{P}}^1}(k)$.
Furthermore, we write
\begin{gather*}
H_X=\pi^\ast H^0({\mathbb{P}}^1,{\mathcal{O}}(1)) \subset H^0(X,{\mathcal{O}}_X(1)).
\end{gather*}

We consider the set ${\mathcal{X}}$ of curves $C$ in $Z$ such that:
\begin{itemize}\itemsep=0pt
\item
$\pi:C\to {\mathbb{P}}^1$ is f\/inite-to-one;
\item
the Douady space ${\mathcal{D}}(Z)$ is smooth at $C$;
\item
the sheaf ${\mathcal{N}}_{C/Z}(-2)$ is acyclic, i.e.\
$h^i({\mathcal{N}}_{C/Z}(-2))=0$ for $i=0,1$.
\end{itemize}
These conditions are open, so ${\mathcal{X}}$ is an open submanifold of ${\mathcal{D}}(Z)$.

A useful characterisation of {\em pure dimension + Cohen--Macaulay + finite cover of ${\mathbb{P}}^1$} is provided by

\begin{Lemma}\label{CM}
Let $\pi:X\to Y$ be a~finite-to-one closed surjective holomorphic map from a~complex space $X$ to a~complex manifold
$Y$.
Then the following conditions are equivalent:
\begin{itemize}\itemsep=0pt
\item[$(i)$] $X$ is pure dimensional and Cohen--Macaulay;
\item[$(ii)$] $\pi$ is flat;
\item[$(iii)$] $\pi_\ast{\mathcal{O}}_X$ is locally free;
\item[$(iv)$] all fibres of $\pi$ have the same length.
\end{itemize}
\end{Lemma}

\begin{proof}
Follows from~\cite[Corollary 18.17]{Eis1} and the equality $\dim_xX=\dim_{\pi(x)}Y+\dim_{x} X_{\pi(x)}$ for f\/lat
maps~\cite[Proposition II.2.11]{GR}.
\end{proof}

Thus, for a~$C\in {\mathcal{X}}$, each f\/ibre $C_\zeta$ is a~$0$-dimensional complex space of constant length $d$, which
we call the {\em degree of the curve $C$}.

For any $C\in {\mathcal{X}}$ and $\zeta\in {\mathbb{P}}^1$, consider the map ${\mathcal{N}}_{C/Z}(-2)\stackrel{\cdot
s}{\longrightarrow} {\mathcal{N}}_{C/Z}(-1)$, where $s\in H_C$ with $s(\zeta)=0$.
This map must be injective, since a~nontrivial kernel sheaf would be supported on a~$0$-dimensional subspace, and so it
would have nontrivial sections, which would then map to nontrivial sections of ${\mathcal{N}}_{C/Z}(-2)$.
We have, therefore, the exact sequence
\begin{gather}
0\to {\mathcal{N}}_{C/Z}(-2)\stackrel{\cdot s}{\longrightarrow} {\mathcal{N}}_{C/Z}(-1)\longrightarrow
{\mathcal{N}}_{C/Z}(-1)|_{C_\zeta}\to 0.
\label{seq}
\end{gather}
Taking cohomology and considering a~generic $\zeta$ shows that
\[
h^0( {\mathcal{N}}_{C/Z}(-1))=dn \qquad \text{and}\qquad h^1(
{\mathcal{N}}_{C/Z}(-1))=0.
\]
Similarly, taking a~section $t$ of ${\mathcal{O}}_{{\mathbb{P}}^1}(2)$ with zeros at $\zeta$ and $\tilde\zeta$ gives the
exact sequence
\begin{gather}
0\to {\mathcal{N}}_{C/Z}(-2)\stackrel{\cdot t}{\longrightarrow} {\mathcal{N}}_{C/Z}\longrightarrow
{\mathcal{N}}_{C/Z}|_{C_\zeta}\oplus {\mathcal{N}}_{C/Z}|_{C_{\tilde\zeta}}\to 0,
\label{seqq}
\end{gather}
from which we derive
\begin{gather}
h^0( {\mathcal{N}}_{C/Z})=2dn,
\qquad
h^1( {\mathcal{N}}_{C/Z})=0.
\label{dim}
\end{gather}
The def\/initions and~\eqref{dim} imply
\begin{Proposition}
The subset ${\mathcal{X}}_d$ of ${\mathcal{X}}$, consisting of curves of degree $d$, is open in ${\mathcal{D}}(Z)$ and
smooth of dimension~$2dn$.
\end{Proposition}

For regularly embedded subspaces some of the assumptions on $C$ are automatically fulf\/illed

\begin{Lemma}\label{ooo}
Let $C$ be a~regularly embedded compact subspace of $Z$ such that $\pi:C\to {\mathbb{P}}^1$ is finite-to-one and
$H^\ast({\mathcal{N}}_{C/Z}(-2))=0$.
Then $C\in {\mathcal{X}}$.
\end{Lemma}

\begin{proof}
The normal sheaf of a~regularly embedded subspace is locally free.
Hence,~\eqref{seq} implies that each f\/ibre $C_\zeta$ has length $h^0( {\mathcal{N}}_{C/Z}(-1))/n$.
Owing to Lemma~\ref{CM}, $C$ is equidimensional and Cohen--Macaulay.
Moreover, for a~regularly embedded $C$,~\eqref{dim} implies that ${\mathcal{D}}(Z)$ is smooth at~$C$.
\end{proof}

Since the linear system $H_C$ is base-free,~\eqref{seqq} implies that the following sequence is also exact:
\begin{gather*}
0\to {\mathcal{N}}_{C/Z}(-2)\to {\mathcal{N}}_{C/Z}(-1)\otimes H_C\to {\mathcal{N}}_{C/Z}\to 0,
%\label{seq-1}
\end{gather*}
and, consequently, there is a~canonical isomorphism
\begin{gather}
H^0({\mathcal{N}}_{C/Z})\simeq H^0({\mathcal{N}}_{C/Z}(-1))\otimes H_C.
\label{tensor}
\end{gather}
We denote by $E$ a~vector bundle over ${\mathcal{X}}_d$, the f\/ibre of which at $C$ is $H^0({\mathcal{N}}_{C/Z}(-1))$.
In the notation of~\eqref{seqq0}, $E= \nu_\ast\bigl({\mathcal{N}}_Y\otimes \mu^\ast{\mathcal{O}}_Z(-1)\bigr)$.
We also write $H$ for the trivial bundle of rank~$2$, with f\/ibre $H_C=\pi^\ast H^0({\mathbb{P}}^1,{\mathcal{O}}(1))$.
The decomposition~\eqref{tensor} induces a~decomposition (of holomorphic vector bundles)
\begin{gather*}
T{\mathcal{X}}_d\simeq E\otimes H.
%\label{decompose}
\end{gather*}
For every $\zeta\in {\mathbb{P}}^1$, we have a~subbundle $Q_\zeta$ of $T{\mathcal{X}}_d$ of rank $dn$ def\/ined as
\begin{gather*}
Q_\zeta=E\otimes s,
\end{gather*}
where $s$ is a~section of ${\mathcal{O}}_{{\mathbb{P}}^1}(1)$ vanishing at $\zeta$.
In other words, the f\/ibre of $Q_\zeta$ at $C$ consists of sections of ${\mathcal{N}}_{C/Z}$ vanishing on $C_\zeta$.

\begin{Proposition}\label{integr}
The distribution $Q_\zeta$ is integrable.
\end{Proposition}

\begin{proof}
Let $X$ and $Y$ be vector f\/ields with values in $Q_\zeta$, and $t\mapsto\gamma_t(m)$, $t\mapsto \delta_t(m)$ their
integral curves beginning at $m$.
The bracket of $[X,Y]$ at any $p\in {\mathcal{X}}_d$ can be computed as $\lim\limits_{t\to 0}\dot\nu(t)/t^2$, where
\begin{gather*}
\nu(t)=\delta_{-t}(\gamma_{-t}(\delta_t(\gamma_t(m)))).
\end{gather*}
Consider the corresponding deformations of curves in $Z$.
Since $X$ and $Y$ take values in $Q_\zeta$, the deformation $\nu(t)$ f\/ixes the f\/ibre $C_\zeta$, and hence $\dot\nu(t)\in
Q_\zeta$.
Thus $[X,Y]_m\in {Q_\zeta}|_m$.
\end{proof}

Now suppose that $Z$ is equipped with an antiholomorphic involution $\sigma$ which covers the antipodal map on
${\mathbb{P}}^1$.
The submanifold $({\mathcal{X}}_d)^\sigma$ of $\sigma$-invariant curves in ${\mathcal{X}}_d$ is either empty or of real
dimension $2dn$.
In the latter case, $({\mathcal{X}}_d)^\sigma$ is canonically a~hypercomplex manifold.
For each complex structure $I_\zeta$, $\zeta\in {\mathbb{P}}^1$, the bundle of $(0,1)$-vectors is $Q_\zeta$.
We state the result as follows.
\begin{Theorem}\label{hc}
Let $Z$ be an equidimensional complex space equipped with a~holomorphic surjection $\pi:Z\to {\mathbb{P}}^1$ and an
antiholomorphic involution $\sigma$ covering the antipodal map on ${\mathbb{P}}^1$.
Then the subset ${\mathcal{M}}_d$ of the smooth locus of ${\mathcal{D}}(Z)$, consisting of $\sigma$-invariant curves $C$
of degree~$d$ such that $H^\ast({\mathcal{N}}_{C/Z}(-2))=0$, is, if nonempty, a~hypercomplex manifold of real dimension
$2d(\dim Z -1)$.
\end{Theorem}
